
Machine learning repositories such as HuggingFace Datasets Hub (60,000+ datasets) lack formal deprecation mechanisms, relying instead on manual README updates that fail to notify users when datasets become stale or are replaced. This gap creates silent failures where models trained on deprecated data degrade without warning. Current practice fails at three infrastructure points: automated candidate detection, task-specific successor mapping, and adoption measurement. This work presents a three-component formal deprecation system integrating automated health metrics (usage velocity, successor emergence, issue signals), context-aware successor graphs (task-conditional replacement paths), and load-time instrumentation (adoption tracking telemetry). Five proof-of-concept experiments (h-e1, h-m1, h-m2, h-m3, h-m4) validate mechanistic feasibility on simulated and mock HuggingFace-compatible data at 100-1000 sample scale. Health metrics achieved 88.3\% precision in deprecation candidate detection (target $\geq 60$\%), context inference achieved 100\% accuracy on synthetic task patterns (target $\geq 70$\%, expected 70-90\% on real data), and load-time instrumentation added 0.21-5\% overhead (target <10\%) with 100\% event capture (95\% CI: 99.05\%-100\%). All gates passed, validating infrastructure mechanisms. Efficacy measurement (adoption lift versus baseline informal mechanisms) was not performed and remains future work. This study demonstrates proof-of-concept infrastructure for dataset lifecycle management, enabling future deployment-scale efficacy testing that neither software package managers nor documentation standards can perform in isolation.
